% !TeX root = ../main.tex
\chapter{多种辐射过程的耦合}
真实情况中的辐射往往是多种辐射相互作用的共同结果。一个简单的物理图像是：不同机制产生种子光子（如韧致辐射和同步辐射），这些光子经过等离子体云的逆康普顿加速，转化为高能电子。在此过程中，我们假设存在一个能够不断被加热以维持自身温度的等离子体云（例如由引力加热）。等离子云会不断制造光子，同时也可以持续与光子散射。这样辐射相互作用的过程就可以简化为一个光子源与恒温等离子体云之间的持续相互作用。
可以简化为：光子源与等离子体云通过康普顿散射相互作用，最终形成高能电子与相应的辐射。我们主要关注韧致辐射、同步辐射和逆康普顿化，而这两者的一个重要区别在于有没有强磁场，所以对上述物理模型的讨论可以分为对有无强磁场情况的讨论。
\n{黑体谱虽然是吸收率为1的极限情况，但是并不是吸收$h\nu$的光子就释放$h\nu$的光子，而更有可能将能量在内部重新分配，释放出多个光子以维系能量守恒。}

\section{轫致自康普顿散射}
热轫致辐射产生的光子频率覆盖广，在恒温且无磁场的非相对论等离子体中，辐射场的演化取决于吸收和散射作用。在不同频段表现出截然不同的物理本征，且最终谱形的呈现依赖于系统的空间尺度。

我们首先假设电子云是足够大的，也就是说有效光深$\tau_{eff}$足够大。这样所有光子都能只考虑散射和吸收作用，而不用考虑逃逸。注意到自由-自由吸收截面$\alpha_{ff} \propto \nu^{-3}$，而散射截面在非相对论极限下维持为汤姆孙截面 $\sigma_T$。在低频能段，光子与电子云的相互作用以吸收为主，光谱表现为黑体谱，或者说它的低频近似——瑞利-金斯谱，这也是一个幂律谱。
\n{我们假设统计上电子热运动的能量是大于光子能量的，只有在充分逆康普顿化后，光子才能接近电子的热能量水平。}
随着频率升高，自由-自由吸收截面快速减弱而散射效应凸显。光子进入散射主导的有效光学厚区域，此时光子在逃逸前因频繁散射而陷入随机行走，路径的显著延长增强了光子被吸收的概率，散射对光子能量的增益使光谱偏离标准黑体谱，在频谱上表现为斜率为 1 的修正黑体谱。

\n{观测到的光子电子云为了平衡内部光子吸收的而射出的光子与内部光子逃逸成分共同构成。}
随着频率进一步升高，虽然单次散射的能量增益仍然有限，但是粒子在被吸收前可以经过更多次数的散射增能，康普顿y参数接近于1。然而，这种硬化存在热力学上限，与上一章机制类似，康普顿参数接近1时，历尽散射后光子群体的能量开始接近电子的热能量，单次散射能量增益近乎于0，大量的光子在多次散射后堆积在小于电子热能的波段，光谱上辐射随着频率以更胜于修正黑体谱的速度增长，直到到达一个峰值，此时电子不具备让光子继续增能的能力，能跨过这一频率的光子少之又少，迅速下降，形成一个明显的峰值。

上述能谱特征是在我们做出了电子云足够大的假设下呈现的，若电子云尺度不足以支撑光子经历充分散射，光子将无法在Wien峰处堆积，退化为轫致辐射谱。
\section{同步自康普顿散射 (SSC)}

当等离子体处于\textbf{非热、有磁场}的环境时，辐射过程转变为 SSC 机制。SSC 描述的是\textbf{同一群电子}的两个连续过程：同步辐射产生种子光子，然后用这些种子光子进行逆康普顿散射。

\begin{itemize}
    \item \textbf{电子能谱假设：} 电子能量遵循幂律分布：
    \begin{equation}
        N(\gamma) = N_0 \gamma^{-p} \quad (\gamma_1 < \gamma < \gamma_2)
    \end{equation}
    \item \textbf{光学薄假设：} 光子经过一次散射后就逃离气体云。
\end{itemize}

\paragraph{为何聚焦于光学薄？}
SSC 理论的简化推导（用于确定谱形 $F_\nu \propto \nu^{-\alpha}$）强烈依赖于以下两个前提：
\begin{enumerate}
    \item \textbf{同步辐射为光学薄 ($\tau_s \ll 1$):} 只有此时，同步辐射的强度才直接反映了电子的真实幂律能谱 $N(\gamma)$。如果同步辐射被吸收（光学厚），其谱形会被温度修正，SSC 的推导基础就会被破坏。
    \item \textbf{SSC 过程（高能部分）:} 虽然 SSC 涉及多次散射的累积，但在计算其强度时，我们利用了\textbf{光学薄近似}来描述种子光子场。此外，高能 SSC 光子自身的散射截面会进入\textbf{克莱因-仁科（KN）区}，导致散射截面减小，从而限制了更高阶散射的发生，使高能谱最终偏离幂律而指数下降。
\end{enumerate}

\subsection{SSC 的辐射谱}

在讨论同步自康普顿散射（SSC）的辐射谱时，其物理图像不同于之前讨论的热平衡气体逆康普顿化。由于 SSC 模型通常建立在\textbf{光学薄} 的假设之上，我们只关注光子在源区内\textbf{仅被散射一次}便逃逸的情况，因此没有必要依据康普顿 $y$ 参数进行复杂的分类讨论。

要计算 SSC 的辐射谱，本质上还是要利用\textbf{幂律分布电子散射通用谱公式：}
\begin{equation}
    P(\nu) = 8\pi r_0^2 ch \iint N(\gamma) f(\nu / 4\gamma^2 \nu_i) n_{ph}(\nu_i) \dd{\nu_i} \dd{\gamma} \propto \nu^{-(p-1)/2}
\end{equation}
\textbf{这个公式的最终确定依赖于入射光源的数密度}。

\paragraph{同步辐射光子场的数密度}
计算 SSC 辐射谱的第一步是确定作为“种子”的同步辐射光子场的数密度。其核心计算思路是：利用已知的同步辐射发射率确定单位体积内的光子产生率，结合光子在球形源内的平均停留时间，求得稳态下的光子数密度分布。

\begin{theorem}[同步辐射场光子数密度]
    \begin{equation}
        \boxed{n_s(\nu) = n_{s0} \nu^{-\frac{p+1}{2}} \approx \frac{3R}{4c} \frac{P_s(\nu)}{h\nu}}
    \end{equation}
\end{theorem}

\begin{proof}
    在半径为 $R$ 的球形源中，光子在源内的平均停留时间由平均逃逸路程决定，估算为 $t_{stay} = \frac{3R}{4c}$。
    设源内同步辐射的单色功率发射率为 $P_s(\nu)$（单位：$\text{erg} \cdot \text{s}^{-1} \cdot \text{Hz}^{-1}$），则单位频率、单位体积的光子产生率为 $\frac{P_s(\nu)}{h\nu V}$。
    在稳态条件下，单位体积内的光子总数密度 $n_s(\nu)$ 为产生率与停留时间的乘积：
    \begin{equation}
        n_s(\nu) = \text{产生率} \times t_{stay} = \frac{P_s(\nu)}{h\nu \cdot V} \cdot \frac{3R}{4c} \cdot V = \frac{3R}{4c} \frac{P_s(\nu)}{h\nu}
    \end{equation}
    已知同步辐射功率 $P_s(\nu) \propto \nu^{-(p-1)/2}$，代入上式：
    \begin{equation}
        n_s(\nu) \propto \frac{\nu^{-\frac{p-1}{2}}}{\nu} = \nu^{-\frac{p+1}{2}}
    \end{equation}
    定义系数 $n_{s0}$ 以包含所有常数项，即得到 $n_s(\nu) = n_{s0} \nu^{-\frac{p+1}{2}}$。
\end{proof}

\paragraph{SSC 辐射功率}
将上述得到的光子数密度代入幂律分布电子的逆康普顿散射谱公式中，通过对种子光子频率进行积分，即可得到最终的 $P_{ssc}$。

\begin{theorem}[SSC 辐射功率公式]
    \begin{equation}
        \boxed{P_{ssc}(\nu) = \pi r_0^2 ch N_0 2^{p+3} \frac{p^2+4p+11}{(p+3)^2(p+1)(p+5)} n_{s0} \ln \Lambda \nu^{-\frac{p-1}{2}}}
    \end{equation}
\end{theorem}

\begin{proof}
    已知幂律分布电子 $N(\gamma) = N_0 \gamma^{-p}$ 对各向同性光子场 $n_{ph}(\nu_i)$ 的单次散射功率公式为：
    \begin{equation}
        P(\nu) = \pi r_0^2 ch N_0 2^{p+3} \frac{p^2+4p+11}{(p+3)^2(p+1)(p+5)} \nu^{-\frac{p-1}{2}} \int \nu_i^{\frac{p-1}{2}} n_{ph}(\nu_i) \dd{\nu_i}
    \end{equation}
    将 SSC 的特定种子场 $n_s(\nu_i) = n_{s0} \nu_i^{-(p+1)/2}$ 代入积分项：
    \begin{align}
        \text{Integral} &= \int_{\nu_1}^{\nu_2} \nu_i^{\frac{p-1}{2}} \cdot \left( n_{s0} \nu_i^{-\frac{p+1}{2}} \right) \dd{\nu_i} \nonumber \\
        &= n_{s0} \int_{\nu_1}^{\nu_2} \nu_i^{\left( \frac{p-1}{2} - \frac{p+1}{2} \right)} \dd{\nu_i} = n_{s0} \int_{\nu_1}^{\nu_2} \nu_i^{-1} \dd{\nu_i} \nonumber \\
        &=  n_{s0} [\ln \nu_i]_{\nu_1}^{\nu_2} = n_{s0} \ln \left( \frac{\nu_2}{\nu_1} \right)
    \end{align}
    定义 $\Lambda = \nu_2 / \nu_1$（同步辐射的频率范围比），则积分项简化为 $n_{s0} \ln \Lambda$。将此积分结果代回原功率公式，即得到上述框选的 $P_{ssc}(\nu)$ 表达式。
\end{proof}

在该公式中，系数 \textbf{$n_{s0} \ln \Lambda$} 集中包含了同步辐射场的所有物理性质。由于 $n_{s0}$ 与同步辐射强度成正比，而同步辐射强烈依赖于磁场强度 $B$，\textbf{通过这个系数，SSC 散射谱完美继承了同步辐射依赖磁场的性质}。此外，结果显示 SSC 谱的光谱指数 $\alpha = (p-1)/2$ 与同步辐射谱完全相同，这意味着在能谱分布上，两个峰的幂律段是相互平行的。

\subsection{SSC 辐射功率与同步辐射功率对比}

为了定量衡量高能辐射中同步自康普顿（SSC）过程的贡献，需要计算 SSC 功率与同步辐射功率的比值。

\begin{theorem}[功率之比]
    \begin{equation}
        \boxed{\frac{P_{ssc}(\nu)}{P_s(\nu)} \approx 9 \sigma_T R N_0 2^{p-2} \ln \Lambda \frac{p^2+4p+11}{(p+3)^2(p+1)(p+5)} \propto \tau_{es} \ln \Lambda}
    \end{equation}
\end{theorem}

\begin{proof}
    在计算 SSC 辐射功率时，首先考虑幂律分布电子对各向同性同步辐射场进行单次逆康普顿散射的通用功率公式，其表达式包含经典电子半径 $r_0$ 构成的系数项：
    \begin{equation}
        P_{ssc}(\nu) = \pi r_0^2 ch N_0 2^{p+3} \frac{p^2+4p+11}{(p+3)^2(p+1)(p+5)} \nu^{-\frac{p-1}{2}} \int_{\nu_1}^{\nu_2} \nu_i^{\frac{p-1}{2}} n_s(\nu_i) \dd{\nu_i}
    \end{equation}
    利用汤姆孙散射截面 $\sigma_T = \frac{8\pi}{3} r_0^2$ 将经典电子半径项替换为 $r_0^2 = \frac{3\sigma_T}{8\pi}$。同时，代入之前确定的同步辐射光子数密度分布 $n_s(\nu_i) = n_{s0} \nu_i^{-(p+1)/2}$，其中系数 $n_{s0} = \frac{3R}{4ch} P_{s0}$。此时，完整的 $P_{ssc}(\nu)$ 表达式转化为包含有效散射光深 $\tau_{es} = N_0 \sigma_T R$ 的形式：
    \begin{equation}
        P_{ssc}(\nu) = \pi \left( \frac{3\sigma_T}{8\pi} \right) ch N_0 2^{p+3} \frac{p^2+4p+11}{(p+3)^2(p+1)(p+5)} \nu^{-\frac{p-1}{2}} \left( \frac{3R}{4ch} P_{s0} \right) \ln \Lambda
    \end{equation}
    在上述表达式中，普朗克常数 $h$ 与光速 $c$ 在分子分母中相互抵消，通过合并项并提取同步辐射的单色功率系数 $P_{s0}$，可以得到 SSC 功率与同步辐射单色功率 $P_s(\nu) = P_{s0} \nu^{-(p-1)/2}$ 之比：
    \begin{equation}
        \frac{P_{ssc}(\nu)}{P_s(\nu)} = \frac{\frac{9}{32} \sigma_T N_0 R 2^{p+3} \frac{p^2+4p+11}{(p+3)^2(p+1)(p+5)} \ln \Lambda P_{s0} \nu^{-\frac{p-1}{2}}}{P_{s0} \nu^{-\frac{p-1}{2}}}
    \end{equation}
    考虑到 $\tau_{es} = N_0 \sigma_T R$ 表示散射光深，整理就得到须证式。
\end{proof}

同步辐射作为一次辐射过程，其强度直接正比于电子数密度 $N_0$；而 SSC 作为一种涉及电子“产生光子”后又“散射光子”的二次过程，其强度应该与电子数密度的平方 $N_0^2$ 成正比，公式可以看出来这一点。这个比值真正依赖的与光子源无关的物理量只要 $\tau_{es}$ 和 $\Lambda$，也就是散射光深和同步辐射的频率范围。

\subsection{SSC 复合光谱图像分析}

典型的 SSC 辐射在能谱分布（SED）上呈现出清晰的双峰结构。左侧低频峰源于电子在磁场中的同步辐射，右侧高频峰则由这些同步光子被同一群电子逆康普顿散射产生。在光学薄的幂律频段，由于两类辐射共享同一套电子能谱分布，SSC 峰完美继承了同步辐射的光谱指数，使得两个辐射峰的幂律段在对数坐标系下表现为相互平行的直线。



% \begin{figure}[H]
%     \centering
%     \includegraphics[width=0.8\textwidth]{Pasted image 20251226211712.png}
%     \caption{SSC 复合能谱图示}
% \end{figure}

由于天体物理习惯使用 $\nu F_\nu$ 作为纵轴，若原能谱流量满足 $F_\nu \propto \nu^{-\alpha}$（其中 $\alpha = \frac{p-1}{2}$），则在能谱图中表现出的视觉斜率为：
\begin{equation}
    \frac{d \log (\nu F_\nu)}{d \log \nu} = 1 - \alpha
\end{equation}
这两个峰之间的垂直高度差（流量比）由有效光深决定，这直观地展示了能量从磁场能向高能光子场转移的效率。而在光谱的极两端，低频的同步辐射受自吸收压制表现为 $\nu^{5/2}$ 的上升，高频的 SSC 辐射则受限于电子加速的最高能量 $\gamma_2$ 以及进入克莱因-仁科区域后的散射截面压制，共同勾勒出高能天体辐射的物理边界。